\documentclass[10pt,reqno]{amsart}
\usepackage{cmap}
\usepackage[T2A]{fontenc}
\usepackage[utf8]{inputenc}
\usepackage[english,ukrainian]{babel}
\usepackage{indentfirst,textcase}
\usepackage[a5paper,margin={15mm,15mm}]{geometry}

\theoremstyle{plain}
\newtheorem{theorem}{Теорема}
\newtheorem*{theorem*}{Теорема}
\newtheorem{corollary}{Наслідок}
\newtheorem*{corollary*}{Наслідок}
\newtheorem{lemma}{Лема}
\newtheorem*{lemma*}{Лема}
\newtheorem{proposition}{Твердження}
\newtheorem*{proposition*}{Твердження}
\newtheorem{conjecture}{Гіпотеза}
\newtheorem*{conjecture*}{Гіпотеза}
\theoremstyle{definition}
\newtheorem{definition}{Означення}
\newtheorem*{definition*}{Означення}
\newtheorem{example}{Приклад}
\newtheorem*{example*}{Приклад}
\theoremstyle{remark}
\newtheorem{remark}{Зауваження}
\newtheorem*{remark*}{Зауваження}

\begin{document}

\title{Глибина Тьюкі деяких багатовимірних імовірнісних розподілів} 

\author{В.~Є.~Панаско}
\address{КПІ ім. Ігоря Сікорського, Київ, Україна}
\email{panasko.vitalii@lll.kpi.ua}

\makeatletter\def\and{\unskip,\;\space\ignorespaces}\def\author@andify{\nxandlist{\and}{\and}{\and}}\makeatother

\maketitle\pagestyle{empty}\thispagestyle{empty}

Функція глибини Тьюкі, вперше запропонована в~\cite{tukey1975}, є однією з найвживаніших мір центральності в багатовимірному статистичному аналізі. Для ймовірнісного розподілу~$\mu$ на~$\mathbb{R}^d$ вона визначається як
\[
D_{\mu}(x) = \inf_{H \ni x} \mu(H),
\]
де інфімум береться по всіх замкнених півпросторах $H \subset \mathbb{R}^d$, що містять точку~$x$. Значення~$D_{\mu}(x)$ характеризує ступінь \lq\lq центральності\rq\rq\ або \lq\lq типовості\rq\rq\ точки $x$ відносно розподілу~$\mu$: ця величина набуває найбільших значень в центральній частині розподілу і спадає до нуля при~$\|x\| \to \infty$. Детальний огляд властивостей функції $D_{\mu}$ наведено в \cite{rousseeuw1999}.

Аналітичне обчислення функції глибини Тьюкі, як правило, є складним і можливе лише для окремих класів розподілів. У двовимірному випадку точні вирази відомі для низки розподілів, зокрема наведених у~\cite{rousseeuw1999}. В нашій роботі ми узагальнюємо ці результати на нові класи багатовимірних розподілів: рівномірний розподіл на кулях у~$\mathbb{R}^d$ (що фактично включає також еліпсоїди завдяки аффінній еквіваріантності функції~$D_{\mu}$) та розподіли з незалежними симетричними~$\alpha$-стійкими координатами. Крім того, для деяких інших двовимірних розподілів --- експоненціального та закону Стьюдента з незалежними координатами --- аналітичний опис глибини є недоступним; у цих випадках ми досліджуємо геометрію відповідних верхніх рівневих множин чисельно та вивчаємо вплив параметрів розподілів на форму контурів.

\begin{thebibliography}{99}
	
	\bibitem{tukey1975}
	Tukey, J.\ W.\ (1975).
	\textit{Mathematics and the picturing of data}.
	In: \textit{Proceedings of the International Congress of Mathematicians} (Vancouver, 1974), vol. 2, pp.~523--531.
	
	\bibitem{rousseeuw1999}
	Rousseeuw, P.\ J., Ruts, I.\ (1999).
	\textit{The depth function of a population distribution}.
	\textit{Metrika}, \textbf{49}, 213--244.
	
\end{thebibliography}
\end{document}
